\begin{textsample}{POS Dim 4 – 2018\_c5 – Score 1.00 – 2018\_c5/enade\_2018\_des\_d0d272ca78.txt}  \label{ex:f4_pos_221}
CAPÍTULO 5 DISTRIBUIÇÃO DOS CONCEITOS 5.1 PANORAMA NACIONAL DA DISTRIBUIÇÃO DOS CONCEITOS O Conceito Enade é um indicador de qualidade que avalia os cursos por intermédio dos desempenhos dos estudantes a partir dos resultados do Enade.
Os valores possíveis vão de 1 ( pior situação ) a 5 ( melhor situação ).
Os cursos com apenas um ou sem nenhum concluinte participante não obtêm o Conceito Enade, ficando Sem Conceito ( SC ).
Na Tabela 5.1, são apresentadas a quantidade e a distribuição de cursos de Design Participantes do Enade/2018, por faixa de conceito e Grande Região.
Enfatiza -se, mais uma vez, que a diferença entre os cursos tabulados neste capítulo e no capítulo 2 corresponde aos cursos Sem Conceito, ou seja àqueles, em princípio, sem alunos concluintes que houvessem participado da prova ou àqueles que tivessem tão somente um aluno concluinte19.
Observando -se os dados da Tabela 5.1, nota -se que, dos 160 cursos Participantes, 56 ( 35,0 \% ) classificaram -se com conceito 3, o valor modal.
Esse foi o conceito modal nas regiões Sudeste ( 43,5 \% ) e Sul ( 34,9 \% ).
Na região Nordeste ( 47,4 \% ), a moda foi o conceito 4 ; já na região Norte ( 60,0 \% ), foi o conceito 2.
Apenas a região Centro-Oeste ( 50,0 \% ) obteve o conceito 5.
Dos 160 cursos de Design, sete ( 4,4 \% ) ficaram Sem Conceito ( SC ).
A região Norte participou com cinco cursos ou 3,1 \% do total nacional.
Desses, três cursos receberam o conceito 2, o que equivale a 60,0 \% do total regional, além de ser o conceito modal na região, como já comentado.
A um curso ( 20,0 \% ) atribuiu -se conceito 4.
Um curso ficou Sem Conceito ( SC ), e nenhum curso recebeu os conceitos 1,3 ou 5 nessa região. 19 Esses cursos com somente um concluinte não foram considerados no capítulo 2.
A região Nordeste participou com 19 cursos ou 11,9 \% do total nacional.
Nessa região, nove cursos ( 47,4 \% em termos regionais ) obtiveram conceito 4, o conceito modal para a região, como já comentado.
Os conceitos 2 e 3 foram atribuídos, respectivamente, a cinco e a três cursos ( 26,3 \% e 15,8 \%, respectivamente ).
Ainda nessa região, os conceitos 1 e 5 foram atribuídos a um curso cada ( 5,3 \% ).
Nenhum curso ficou Sem Conceito ( SC ).
Dos 69 cursos Participantes da região Sudeste ( 43,1 \% do total nacional ), o conceito 3 foi o conceito modal, como já anteriormente assinalado, tendo sido obtido por 30 cursos ( 43,5 \% ).
O conceito 1 foi atribuído a dois cursos ( 2,9 \% ), e o conceito 2, a 13 cursos ( 18,8 \% ).
Receberam os conceitos 4 e 5, respectivamente, 19 cursos ( 27,5 \% ) e três cursos ( 4,3 \% ).
Nessa região, dois cursos ficaram Sem Conceito ( SC ).
Os 63 cursos da região Sul corresponderam a 39,4 \% do total nacional.
Como já indicado, houve predominância do conceito 3, atribuído a 34,9 \% dos cursos da região ( conceito modal ), ou seja, atribuído a 22 dos 63 cursos Participantes na região Sul.
O conceito 1 foi atribuído a dois cursos ( 3,2 \% ), e o conceito 2, a dez cursos ( 15,9 \% ).
Os conceitos 4 e 5, foram atribuídos a, respectivamente, 21 cursos ( 33,3 \% ) e quatro cursos ( 6,3 \% ).
Quatro cursos ficaram Sem Conceito ( SC ) na região Sul.
Já dos quatro cursos Participantes na região Centro-Oeste ( 2,5 \% do total nacional ), dois ( 50,0 \% em termos regionais ) receberam conceito 5, o conceito modal, como já fora destacado.
Os conceitos 2 e 3 foram atribuídos a um curso cada ( 25,0 \% ).
Nenhum curso ficou Sem Conceito ( SC ), assim como também nenhum curso recebeu os conceitos 1 ou 4.
No Gráfico 5.1, é apresentada a distribuição cumulativa do Conceito Enade, segundo Grande Região.
Em um gráfico de distribuição cumulativa, em cada valor de abscissa ( eixo x ), o valor apresentado no eixo das ordenadas ( eixo y ) é o valor acumulado da variável.
Por exemplo, a região Norte ( linha azul ) 75,0 \% dos cursos com conceito 3 ou abaixo desse valor.
A região Centro-Oeste ( linha laranja ), por outro lado, apresenta 50,0 \% dos cursos com conceito 3 ou abaixo.
Em linhas gerais, em um gráfico de distribuição cumulativa, poligonais mais à esquerda representam Grandes Regiões com pior distribuição, e poligonais mais à direita, Grandes Regiões com melhores conceitos.
A região Norte ( linha azul ) apresenta os cursos com conceitos menores, ao passo que a região Centro-Oeste ( linha laranja ) apresenta, grosso modo, a distribuição com conceitos maiores.
Cursos em IES situadas nas regiões Sudeste ( linha vermelha ), Nordeste ( linha verde ) e Sul ( linha rosa ) apresentam uma situação intermediária. 5.2 CONCEITOS POR CATEGORIA ADMINISTRATIVA, POR MODALIDADE DE ENSINO E GRANDE REGIÃO Na Tabela 5.2, é apresentada a distribuição dos cursos de Design Participantes do Enade/2018, por Categoria Administrativa e Modalidade de Ensino, de acordo com os conceitos por eles alcançados, segundo as Grandes Regiões brasileiras.
Dos 160 cursos Participantes, 49 ( 30,6 \% ) eram ministrados em Instituições Públicas, e 111 ( 69,4 \% ), em Instituições Privadas.
Quanto à Modalidade de Ensino, todos os 160 cursos eram Presenciais.
De acordo com as informações da Tabela 5.2, em termos nacionais, entre Instituições Públicas e Privadas, observa -se que, dos dez cursos avaliados com conceito 5, todos eram oferecidos em IES Públicas.
Dos 49 cursos Participantes de IES Públicas, o conceito 4 foi o valor modal, atribuído a 27 cursos ( 55,1 \% ).
Dos demais cursos de Instituições Públicas Participantes, um curso ( 2,0 \% ) foi avaliado com conceito 1, três cursos ( 6,1 \% ) foram avaliados com conceito 2, seis ( 12,2 \% ) cursos, com conceito 3 e, como já comentado, dez cursos ( 20,4 \% ), com conceito 5.
Dois cursos ficaram Sem Conceito ( SC ).
Na privada, o conceito modal foi 3, atribuído a 50 cursos, o correspondente a 45,0 \% dos 111 cursos da categoria.
Dos demais cursos Participantes, quatro ( 3,6 \% ) receberam conceito 1, e 29 ( 26,1 \% ), conceito 2.
O conceito 4 foi atribuído a 23 cursos ( 20,7 \% ).
Nessa Categoria Administrativa, cinco cursos ( 4,5 \% ) ficaram Sem Conceito ( SC ), e nenhum curso obteve conceito 5.
Dos 160 cursos Participantes da modalidade Presencial, o conceito 3 foi o valor modal, atribuído a 56 cursos ( 35,0 \% ).
Dos demais cursos nessa Modalidade de Ensino, cinco cursos ( 3,1 \% ) receberam conceito 1, 32 cursos ( 20,0 \% ) receberam conceito 2, 50 cursos ( 31,3 \% ), conceito 4, e dez cursos ( 6,3 \% ), conceito 5.
Nesta Modalidade de Ensino, sete cursos ( 4,4 \% ) ficaram Sem Conceito ( SC ).
Na análise por região, observa -se que, na região Norte, as Instituições Privadas participaram com quatro cursos ( 80,0 \% do total regional, a maior proporção entre as regiões ), dos quais um ficou Sem Conceito ( SC ), e nenhum recebeu os conceitos 1, 3, 4 ou 5.
O conceito modal para as Instituições Privadas na região foi 2, atribuído a três cursos, correspondendo a 75,0 \% dessa categoria na região.
As Instituições Públicas participaram com um curso na região Norte ( 20,0 \% do total regional ), ao qual foi atribuído o conceito 4.
Nenhum curso ficou Sem Conceito ( SC ) ou com conceito 1, 2, 3 ou 5 nessa combinação de categoria e região.
Todos os cursos oferecidos na região Norte eram presenciais.
Na região Nordeste, a rede privada concentrou 8 dos 19 cursos Participantes, o equivalente a 42,1 \% do total da região.
Dentre os cursos oferecidos por IES Privadas no Nordeste, quatro cursos foram avaliados com conceito 2, o conceito modal.
Dois cursos obtiveram conceito 3.
Os conceitos 1 e 4 foram atribuídos a um curso, cada.
Nessa combinação de categoria e região, nenhum curso ficou Sem Conceito ( SC ) e nenhum recebeu conceito 5.
As Instituições Públicas da região Nordeste participaram com 11 cursos ( 57,9 \% do total da região ).
Desses, oito obtiveram conceito 4, o valor modal.
Os conceitos 2, 3 e 5 foram atribuídos a um curso, cada.
Nenhum curso ficou Sem Conceito ( SC ) ou obteve conceito 1 nessa região.
Todos os cursos oferecidos na região Nordeste também eram presenciais.
Na região Sudeste, a proporção de cursos da rede privada, 76,8 \%, foi a segunda mais elevadas dentre as regiões brasileiras, correspondendo a 53 dos 69 cursos Participantes.
Nessa categoria e região, o conceito modal foi 3, atribuído a 27 cursos.
Dois cursos ficaram Sem Conceito ( SC ).
Os demais foram avaliados com conceito 1 ( um curso ), conceito 2 ( 13 cursos ) e conceito 4 ( dez cursos ).
Dos 16 cursos oferecidos por Instituições Públicas na região Sudeste, o conceito modal foi 4, atribuído a nove cursos.
Nessa região, nenhum curso de IES Públicas ficou Sem Conceito ( SC ).
Um curso recebeu conceito 1 e foram atribuídos conceitos 3 e 5 a três cursos cada.
Também na região Sudeste, todos os 69 cursos foram oferecidos na modalidade Presencial.
As Instituições Privadas concentraram 44 dos 63 cursos Participantes da região Sul, 69,8 \% do total regional.
Desses, 22 cursos obtiveram conceito 3, o conceito modal.
Nesta combinação de Categoria Administrativa e Grande Região, oito cursos receberam conceito 2 e outros 12, conceito 4.
Dois cursos receberam conceito 1, e outros dois cursos ficaram Sem Conceito ( SC ).
Nenhum curso recebeu conceito 5.
As Instituições Públicas na região Sul participaram com 19 cursos ( 30,2 \% dos cursos da região ), aos quais foram atribuídos os conceitos 2 e 3 ( dois cursos, cada ), 4 ( nove cursos, conceito modal ) e 5 ( quatro cursos ).
Dois cursos da região Sul oferecidos por IES Pública ficaram Sem Conceito ( SC ).
Nenhum curso recebeu conceito 1.
Todos os cursos oferecidos na região Sul eram presenciais.
Na região Centro-Oeste, dois dos quatro cursos Participantes eram de Instituições Privadas ( 50,0 \% em termos regionais ).
Desses, um curso ( 50,0 \% ) recebeu conceito 3, juntamente com o conceito 2 ( também um curso, 50,0 \% ).
Nenhum curso ficou Sem Conceito ( SC ), e nenhum curso recebeu os conceitos 1, 4 ou 5.
Quanto aos dois cursos oferecidos por Instituições Públicas na região Centro-Oeste ( 50,0 \% do total regional ), ambos apresentaram conceito 5, conceito modal.
Nenhum curso ficou Sem Conceito ( SC ) ou recebeu os demais conceitos nessa combinação de categoria e região.
Todos os cursos oferecidos na região Centro-Oeste também eram presenciais.
No Gráfico 5.2, apresenta -se a distribuição do Conceito Enade, segundo a Categoria Administrativa da IES.
Os cursos em IES Públicas ( linha azul ) apresentam uma poligonal mais à direita do que os em IES Privadas ( linha verde ) e, consequentemente, uma distribuição de Conceitos Enade com valores maiores. 5.3 CONCEITOS POR ORGANIZAÇÃO ACADÊMICA E GRANDE REGIÃO Na Tabela 5.3, encontra -se a distribuição dos conceitos atribuídos aos cursos Participantes do Enade/2018, na Área de Design, por Organização Acadêmica, segundo as Grandes Regiões brasileiras.
Dos 160 cursos de Design Participantes, 104 eram oferecidos em Universidades ; 24, em Centros Universitários ; 31, em Faculdades ; e um, em CEFET/IFET.
Esta distribuição corresponde a, respectivamente, 65,0 \%, 15,0 \%, 19,4 \% e 0,6 \% dos cursos.
De acordo com os dados apresentados, de todos os dez cursos avaliados com conceito 5, todos eram vinculados a Universidades.
Os cursos vinculados a esse tipo de Organização Acadêmica tiveram o conceito 4 como conceito modal, atribuído a 44 cursos ( 42,3 \% ).
Os demais cursos vinculados a Universidades receberam os conceitos 1 ( um curso ), 2 ( 14 cursos ), 3 ( 31 cursos ) e 5 ( dez cursos, como já mencionado ).
Ficaram Sem Conceito ( SC ) quatro cursos.
Dos cursos em Centros Universitários, o conceito modal foi 3, atribuído a 15 cursos ( 62,5 \% ).
Apenas um curso vinculado a esse tipo de Organização Acadêmica ficou Sem Conceito ( SC ).
Os outros cursos receberam os conceitos 1 ( um curso ), 2 ( cinco cursos ) e 4 ( dois cursos ).
Nenhum curso obteve conceito 5.
Dos 31 cursos mantidos por Faculdades, 13 ( 41,9 \% ) receberam conceito 2, o conceito modal.
O segundo conceito mais frequentemente atribuído aos cursos oferecidos por essa categoria de IES foi o conceito 3, recebido por nove cursos.
Na sequência, três cursos obtiveram conceito 1 e quatro cursos, o conceito 4.
Dos cursos oferecidos em Faculdades, dois ficaram Sem Conceito ( SC ) e nenhum obteve conceito 5.
O curso mantido por CEFET/IFET recebeu conceito 3.
Nenhum curso ficou Sem Conceito ( SC ) ou recebeu os demais conceitos nessa combinação de categoria e região.
Considerando -se separadamente as regiões brasileiras, verifica -se que, na região Norte, as Universidades ( 20,0 \% ) participaram com apenas um dos cinco cursos avaliados, ao ao qual foi atribuído conceito 4.
Nenhum curso ficou Sem Conceito ( SC ) ou recebeu os demais conceitos nessa combinação de categoria e região.
Nenhum curso de Centros Universitários participou da avaliação na região Norte.
As Faculdades ( 80,0 \% ) foram representadas por quatro cursos, três dos quais receberam conceito 2, o conceito modal.
Apenas um curso ficou Sem Conceito ( SC ).
Nenhum curso recebeu os conceitos 1,3,4 ou 5.
Nenhum curso oferecido por CEFET/IFET participou na região Norte.
Na região Nordeste, as Universidades participaram com 14 dos 19 cursos da Área de Design oferecidos.
O conceito modal 4 foi atribuído a nove cursos ( 64,3 \% ).
O conceito 2 foi atribuído a três cursos e os conceitos 3 e 5 foram atribuídos a um curso, cada.
Nenhum curso oferecido por Universidades no Nordeste recebeu conceito 1 ou ficou Sem Conceito ( SC ).
Os Centros Universitários contaram com três cursos Participantes na região Nordeste, dois dos quais receberam o conceito modal 2.
Apenas um curso recebeu conceito 3.
Nenhum curso oferecido por Centros Universitários no Nordeste ficou Sem Conceito ( SC ) ou recebeu os conceitos 1, 4 ou 5.
As Faculdades foram representadas por dois cursos na região Nordeste.
Desses, um curso ( 50,0 \% ) recebeu conceito 3, o conceito modal, juntamente com o conceito 1 ( também um curso e 50,0 \% ).
Na região Nordeste, nenhum curso mantido por Faculdades ficou Sem Conceito ( SC ) ou obteve os conceitos 2, 4 e 5.
Nenhum curso oferecido por CEFET/IFET participou na região Nordeste.
Na região Sudeste, as Universidades concentraram 38 dos 69 cursos de Design da região.
Dos cursos oferecidos em Universidades, nessa região, o conceito modal foi 4, atribuído a 16 cursos ( 42,1 \% ).
Os demais cursos receberam os conceitos 1 ( um curso ), 2 ( quatro cursos ), 3 ( 13 cursos ) e 5 ( três cursos ).
Um curso ficou Sem Conceito ( SC ).
Os Centros Universitários participaram com 14 cursos na região Sudeste, dos quais 11 obtiveram o conceito modal 3.
O conceito 1 foi atribuído a um curso, e o conceito 2, a dois cursos.
Nenhum curso ficou Sem Conceito ( SC ) nessa combinação de organização acadêmica e região, e nenhum curso obteve os conceitos 4 e 5.
Já as Faculdades foram representadas por 17 cursos na região Sudeste, dos quais um ficou Sem Conceito ( SC ).
O conceito modal 2 foi atribuído a sete cursos.
Aos demais cursos foram atribuídos os conceitos 3 ( seis cursos ) e 4 ( três cursos ).
Nenhum curso obteve conceito 1 ou 5.
Nenhum curso oferecido por CEFET/IFET participou na região Sudeste.
Dos 63 cursos da região Sul, 47 eram mantidos por Universidades, para os quais o conceito modal foi 4, atribuído a 18 cursos ( 38,3 \% ).
O conceito 2 foi atribuído a seis cursos, o conceito 3 a 16 cursos, e o conceito 5, a quatro cursos.
Nesse tipo de organização, três cursos ficaram Sem Conceito ( SC ), e nenhum recebeu conceito 1.
Dos sete cursos Participantes de Centros Universitários da região Sul, somente um ficou Sem Conceito ( SC ), enquanto nenhum recebeu os conceitos 1 ou 5.
Três cursos receberam o conceito modal 3.
Os conceitos restantes, 2 e 4, foram atribuídos a um curso e a dois cursos, respectivamente.
Quanto aos oito cursos vinculados a Faculdades na região Sul, três receberam o conceito modal 2.
Os conceitos 1 e 3 foram atribuídos a dois cursos, cada e o conceito 4, a um curso.
Nenhum curso ficou Sem Conceito ( SC ) ou recebeu conceito 5.
Já o CEFET/IFET participou com apenas um curso na região Sul, o qual recebeu conceito 3.
Apenas os cursos oferecidos por Universidades participaram da avaliação na região Centro-Oeste.
Para esse tipo de Organização Acadêmica, dois cursos ficaram com o conceito modal 5.
Os conceitos 2 e 3 foram atribuídos a um curso, cada.
Nenhum curso ficou Sem Conceito ( SC ) ou recebeu os conceitos 1 ou 4.
No Gráfico 5.3, apresenta -se a distribuição do Conceito Enade segundo a Organização Acadêmica da IES.
Foi avaliado somente um curso em CEFET/IFET ( linha laranja ) e a poligonal corresponde a um único degrau e, por isso, não entrará na análise que se segue.
Os cursos em Universidades ( linha azul ) apresentam uma poligonal mais à direita do que as demais, denotando uma melhor distribuição de Conceitos Enade.
No outro extremo, com distribuição muito semelhante entre si, encontram -se os cursos mantidos por Centros Universitários ( linha verde ) e Faculdades ( linha vermelha ), apresentando as respectivas poligonais mais à esquerda e, consequentemente, uma distribuição de Conceitos Enade com menores valores.

% matched lemmas: 
\end{textsample}
